\section*{ID9998: Violation of Flat Universe Cosmology: Ω\_rot > 0}
\paragraph{Metadata.}
PiID: 7856672279661 | RTEID: RTE:P38824.4543:T1.1389 | UUID: fcdb5104-6511-5b77-aaea-48e4b2457ed2

\paragraph{Canonical Statement.}
cannot vanish for rotating systems. The constraint \$\textbackslash{}Omega\_\{\textbackslash{}text\{total\}\} = 1\$ required by observations combined with \$\textbackslash{}Omega\_\{\textbackslash{}text\{rot\}\} > 0\$ implies \$\textbackslash{}Omega\_k \textbackslash{}neq 0\$, establishing non-zero curvature.

\paragraph{Canonical Equation.}
\[
\Omega_{\text{rot}} > 0
\]

\paragraph{Dictionary.}
Violation of Flat Universe Cosmology: Ω\_rot > 0 — Ω\_rot > 0

\paragraph{Encyclopedia.}
Violation of Flat Universe Cosmology: Ω\_rot > 0 is an inequality / bound, written Ω\_rot > 0. It states an inequality or bound that constrains the admissible range of the quantity. Within the Trawin Zero Point registry it serves as one element of the framework's mathematical narrative.
